\section{Provider-arkitekturen}
\label{sec:providers}

Kjernen i plugin-arkitekturen er grensesnittet \texttt{IVideoProvider},
som definerer åtte asynkrone metoder for romoperasjoner og
strømstyring. Ingen metoder har default-implementasjon. En provider som
ikke støtter en operasjon kaster \texttt{NotSupportedException}, som
\texttt{GlobalExceptionHandler} oversetter til HTTP~501. Provider velges
ved oppslag i ASP.NET Cores keyed DI-container via
\texttt{VideoProviderFactory}, som normaliserer nøkkelen til lowercase og
kaster \texttt{ArgumentException} ved ukjent provider.
Listing~\ref{lst:ivideoproivder} viser grensesnittet.

\begin{lstlisting}[language={[Sharp]C}, label=lst:ivideoproivder,
  caption={IVideoProvider-grensesnittet.}]
public interface IVideoProvider
{
    Task<VideoRoom> CreateRoomAsync(CreateRoomRequest request,
        CancellationToken ct = default);
    Task<VideoRoom> GetRoomAsync(string roomId,
        CancellationToken ct = default);
    Task DeleteRoomAsync(string roomId, CancellationToken ct = default);
    Task StartStreamAsync(string roomId, CancellationToken ct = default);
    Task StopStreamAsync(string roomId, CancellationToken ct = default);
    Task SetRoomTokenAsync(string roomId, string token,
        CancellationToken ct = default);
    Task<WebRtcLinkResponse> GenerateWebRtcTokenAsync(string roomId,
        string type, CancellationToken ct = default);
    Task<bool> IsAvailableAsync(CancellationToken ct = default);
}
\end{lstlisting}

Begge providers registreres med en typed \texttt{HttpClient} og en
\texttt{DelegatingHandler} for autentisering bakt inn i HTTP-pipelinen,
som vist i listing~\ref{lst:providerregistrering}.

\begin{lstlisting}[language={[Sharp]C}, label=lst:providerregistrering,
  caption={Registrering av providers som keyed services.}]
services.AddTransient<AntMediaAuthHandler>();
services.AddHttpClient<AntMediaProvider>(client => {
    client.BaseAddress = new Uri(baseUrl.TrimEnd('/') + "/");
    client.Timeout = TimeSpan.FromSeconds(10);
}).AddHttpMessageHandler<AntMediaAuthHandler>();
services.AddKeyedTransient<IVideoProvider>("antmedia",
    (sp, _) => sp.GetRequiredService<AntMediaProvider>());

services.AddSingleton<NHNTokenCache>();
services.AddTransient<NHNAuthHandler>();
services.AddHttpClient<NHNVideoProvider>(client => {
    client.BaseAddress = new Uri(baseUrl.TrimEnd('/') + "/");
    client.Timeout = TimeSpan.FromSeconds(10);
}).AddHttpMessageHandler<NHNAuthHandler>();
services.AddKeyedTransient<IVideoProvider>("nhn",
    (sp, _) => sp.GetRequiredService<NHNVideoProvider>());
\end{lstlisting}

\texttt{AntMediaProvider} er strømorientert. Kjerneabstraksjonen i
AntMedia er en \textit{broadcast} med tilstandene \texttt{created},
\texttt{broadcasting} og \texttt{finished}. Start og stopp er eksplisitte
API-kall, og aksess skjer via kortlevde WebRTC-tokens med 90 sekunders
levetid. \texttt{AntMediaAuthHandler} signerer JWT-tokens lokalt med
HMAC-SHA256 og cacher dem i fem minutter med trådsikker refresh via
\texttt{SemaphoreSlim}.

\texttt{NHNVideoProvider} er møteorientert. Et NHN-møte er en booking
med start- og sluttid, og tilgang skjer via en Pexip-basert
konferanse-URL med innebygd PIN som NHN genererer ved opprettelse. Siden
møtet er aktivt så lenge bookingen eksisterer, finnes det ingen
start- eller stopp-operasjon, og \texttt{SetRoomTokenAsync} er heller
ikke aktuell. Disse tre metodene kaster \texttt{NotSupportedException},
som listing~\ref{lst:nhnsupported} viser.

\begin{lstlisting}[language={[Sharp]C}, label=lst:nhnsupported,
  caption={NHN-operasjoner som ikke er støttet.}]
public Task StartStreamAsync(string roomId, CancellationToken ct = default)
    => throw new NotSupportedException(
        "Start stream is not supported for NHN-rom.");

public Task StopStreamAsync(string roomId, CancellationToken ct = default)
    => throw new NotSupportedException(
        "Stop stream is not supported for NHN-rom.");
\end{lstlisting}

Autentisering mot NHN skjer via OAuth~2.0 password grant, og tokenet
caches i singletonen \texttt{NHNTokenCache} i inntil 13 dager.
Kontroller-laget ser ingen av disse forskjellene. En ny videomotor kan
legges til ved å implementere \texttt{IVideoProvider} og registrere den
med en ny nøkkel i \texttt{ProviderExtensions}.